% !TEX TS-program = lualatex
% !TEX encoding = UTF-8 Unicode
\documentclass{article}
\usepackage{lmodern}
\usepackage{graphicx}
\graphicspath{{doc/assets/}{../doc/assets/}{assets/}{../assets/}{../}}

\usepackage[
  paperwidth=338mm, paperheight=190mm,
  top=0mm, bottom=0mm, left=0mm, right=0mm
]{geometry}

% -- ROBUST FONT FIX FOR LUALATEX & OVERLEAF --
\usepackage{fontspec}
\IfFontExistsTF{Roboto}{
  \setsansfont{Roboto}
}{
  \IfFontExistsTF{Arial}{
    \setsansfont{Arial}[
      ItalicFont={Arial},
      ItalicFeatures={FakeSlant=0.25},
      BoldFont={Arial Bold},
      BoldItalicFont={Arial Bold},
      BoldItalicFeatures={FakeSlant=0.25}
    ]
  }{
    \IfFontExistsTF{Noto Sans}{
      \setsansfont{Noto Sans}
    }{
      \IfFontExistsTF{Liberation Sans}{
        \setsansfont{Liberation Sans}
      }{
        \IfFontExistsTF{TeX Gyre Heros}{
          \setsansfont{TeX Gyre Heros}
        }{
          % Fallback mặc định
        }
      }
    }
  }
}
\renewcommand{\familydefault}{\sfdefault}
% ---------------------------------------------

\usepackage{xcolor}
\usepackage{tikz}
\usepackage{fontawesome5}
\usepackage{pgfplots}
\pgfplotsset{compat=1.18}
\usetikzlibrary{calc, matrix, arrows.meta, shapes.geometric}
\usepackage[hidelinks]{hyperref}
\pagestyle{empty}
\parindent=0pt

\newcommand{\TOCitem}[1]{{\color{sap}\vrule width 2pt height 8pt depth 0pt}\ #1}

%── Bảng màu chuẩn hóa thanh lịch (Pure Blue/Yellow)
\definecolor{sap}{HTML}{0000FF}           % Primary: Xanh lam thuần
\definecolor{cyel}{HTML}{FFFF00}         % Secondary: Vàng thuần
\definecolor{cslate}{HTML}{64748B}
\definecolor{cneg}{HTML}{D32F2F}
\definecolor{cpos}{HTML}{2E7D32}
\definecolor{cgry}{HTML}{757575}
\definecolor{clightgray}{HTML}{F8F9FA}

% Màu nền khối Insight Box
\definecolor{cgr}{HTML}{2E7D32}
\definecolor{cbl}{HTML}{1565C0}
\definecolor{cbk}{HTML}{424242}
\colorlet{bggr}{cgr!8}
\colorlet{bgbl}{cbl!8}
\colorlet{bgbk}{cbk!8}

\newcommand{\TitleLine}[1]{%
  \draw[cyel, line width=6pt] (17mm, -10mm) -- (17mm, -25mm);
  \node[anchor=west, font=\fontsize{35}{40}\selectfont\bfseries, text=sap] at (21mm, -17.5mm) {#1};
}

\begin{document}

% ===================================================================
% SLIDE 10: Sơ đồ gameLoop (Kế thừa dữ liệu logic hộp xử lý từ PDF Trang 1)
% ===================================================================
\begin{tikzpicture}[remember picture, overlay]
  \coordinate (NW) at (current page.north west);
  \begin{scope}[shift={(NW)}]
    \TitleLine{Sơ đồ gameLoop}

    \begin{scope}[
        shift={(170mm, -105mm)},
        scale=1.1,
        transform shape,
        box/.style={
            rectangle, draw=black!80, thick, text width=8cm, align=left, rounded corners=3mm, inner sep=10pt, font=\small
        },
        arrow/.style={
            -{Stealth[scale=1.5]}, line width=1.5pt, draw=gray!80
        },
        label/.style={
            font=\footnotesize\bfseries, align=center, text=black!70
        }
    ]

        \node[box, fill=blue!5] (story) at (0, 4.5) {
            \textbf{\large 1. MỞ KHOÁ TUYẾN TRUYỆN}\\
            \textit{(Phân hệ cốt truyện)}\\[1ex]
            $\bullet$ \textbf{Đầu vào:} Điểm và thời gian của vòng chơi trước.\\
            $\bullet$ \textbf{Tác động:} Phân tích di sản để bẻ nhánh cốt truyện. Tạo lựa chọn mới với độ khó tương ứng.\\
            $\bullet$ \textbf{Đầu ra:} Bộ thử thách ẩn để mở nhánh mới.
        };

        \node[box, fill=purple!5] (console) at (5, -3) {
            \textbf{\large 2. THỬ THÁCH TUYẾN TRUYỆN}\\
            \textit{(Phân hệ bảng điều khiển)}\\[1ex]
            $\bullet$ \textbf{Đầu vào:} Bộ thử thách và thông số từ phân hệ cốt truyện.\\
            $\bullet$ \textbf{Tác động:} Cập nhật giao diện. Hiển thị chủ đề tương ứng và mục tiêu.\\
            $\bullet$ \textbf{Đầu ra:} Bộ thông số cấu hình truyền thẳng vào phân hệ lõi trò chơi.
        };

        \node[box, fill=green!5] (core) at (-5, -3) {
            \textbf{\large 3. VƯỢT QUA THỬ THÁCH}\\
            \textit{(Phân hệ lõi trò chơi)}\\[1ex]
            $\bullet$ \textbf{Đầu vào:} Cấu hình độ khó, tốc độ rơi, điều kiện mở khóa.\\
            $\bullet$ \textbf{Tác động:} Xử lý logic giải đố. Người chơi vượt thử thách để lấy quyền năng, đẩy cao thành tích.\\
            $\bullet$ \textbf{Đầu ra:} Kết quả thực chiến lưu lại làm dữ liệu đầu vào.
        };

        \draw[arrow] (story.south east) to[bend left=20] node[label, right=2mm] {Chuyển giao\\Thử thách \& Cấu hình} (console.north);
        \draw[arrow] (console.west) to[bend left=15] node[label, below=2mm] {Kích hoạt\\lối chơi cốt lõi} (core.east);
        \draw[arrow, draw=cneg] (core.north) to[bend left=20] node[label, left=2mm, text=cneg] {\textbf{TÁI KHỞI ĐỘNG}\\Trả kết quả thực chiến\\để định đoạt tương lai} (story.south west);
    \end{scope}
  \end{scope}
\end{tikzpicture}
\newpage

% ===================================================================
% SLIDE 11: Chiến lược công nghệ
% ===================================================================
\begin{tikzpicture}[remember picture, overlay]
  \coordinate (NW) at (current page.north west);
  \begin{scope}[shift={(NW)}]
    \TitleLine{Chiến lược công nghệ}

    \matrix (techtable) [
      matrix of nodes,
      nodes in empty cells,
      nodes={anchor=center, font=\fontsize{16}{22}\selectfont\color{cbk}, text width=65mm, align=center, inner ysep=15pt},
      row 1/.style={nodes={font=\fontsize{20}{26}\selectfont\bfseries\color{sap}, align=center}},
      column 1/.style={nodes={align=left, text width=75mm, font=\fontsize{20}{24}\selectfont\bfseries\color{cslate}}},
      anchor=north west
    ] at (17mm, -55mm) {
      ~ & V1 & V2 & V3 \\

      SDL3 &
      { \Large\faImage \\[1ex] Media } &
      { \Large\faVolumeUp \\[1ex] Sound } &
      { \Large\faGlobe \\[1ex] CDN } \\

      SQLite &
      ~ &
      { \Large\faDatabase \\[1ex] Storage } &
      { \Large\faSync \\[1ex] Sync } \\

      Cloudflare D1 & 
      ~ & 
      ~ & 
      { \Large\faPlug \\[1ex] API } \\
    };

    \draw[cyel, line width=4pt] ([yshift=-2mm]techtable-1-1.south west) -- ([yshift=-2mm]techtable-1-4.south east);
    \draw[clightgray!80!cslate, dashed, line width=1pt] ([yshift=-2mm]techtable-2-1.south west) -- ([yshift=-2mm]techtable-2-4.south east);
    \draw[clightgray!80!cslate, dashed, line width=1pt] ([yshift=-2mm]techtable-3-1.south west) -- ([yshift=-2mm]techtable-3-4.south east);
  \end{scope}
\end{tikzpicture}
\newpage

% ===================================================================
% SLIDE 12: Sơ đồ lớp và CSDL (Kế thừa cấu trúc phân lớp Class từ PDF Trang 6)
% ===================================================================
\begin{tikzpicture}[remember picture, overlay]
  \coordinate (NW) at (current page.north west);
  \begin{scope}[shift={(NW)}]
    \TitleLine{Sơ đồ lớp và CSDL}

    \begin{scope}[
        shift={(17mm, -45mm)},
        classbox/.style={rectangle, draw=cslate, line width=1.5pt, rounded corners=4pt, fill=clightgray, font=\fontsize{12}{16}\selectfont\color{cbk}, text width=55mm, align=center, inner sep=8pt},
        layerlabel/.style={font=\fontsize{16}{20}\selectfont\bfseries\color{sap}, anchor=west},
        arrow/.style={->, line width=1.5pt, draw=cslate}
    ]
        % ================= LAYER 1: gameStory =================
        \node[layerlabel] at (0, 0) {1. Phân hệ Tuyến truyện};
        \node[classbox] (story_db) at (115mm, 0) {\textbf{gameStory\_db.h}\\(SQLite \& Gist API Sync)};
        \node[classbox] (story_layout) at (190mm, 0) {\textbf{gameStory\_layout.h}\\(States \& Transitions)};
        \node[classbox] (story_assets) at (265mm, 0) {\textbf{gameStory\_assets}\\(SVG Logos \& TTF Fonts)};
        \draw[arrow] (story_db) -- (story_layout);
        \draw[arrow] (story_assets) -- (story_layout);

        \draw[cslate, dashed, line width=1pt] (0, -20mm) -- (295mm, -20mm);

        % ================= LAYER 2: gameConsole =================
        \node[layerlabel] at (0, -40mm) {2. Phân hệ Cấu hình};
        \node[classbox] (console_layout) at (115mm, -40mm) {\textbf{gameConsole\_layout.h}\\\texttt{SettingsConfig} Struct};
        \node[classbox] (console_db) at (190mm, -40mm) {\textbf{gameConsole\_db.h}\\\texttt{StoryRow} \& \texttt{Records}};
        \node[classbox] (console_sort) at (265mm, -40mm) {\textbf{gameConsole\_sort.h}\\Smart Sorting Engine (7 Algos)};

        \draw[<->, line width=1.5pt, draw=cslate] (console_layout) -- (console_db);
        \draw[arrow] (console_sort) -- (console_db);

        \draw[cslate, dashed, line width=1pt] (0, -60mm) -- (295mm, -60mm);

        % ================= LAYER 3: gameCore =================
        \node[layerlabel] at (0, -80mm) {3. Phân hệ Lõi trò chơi};
        \node[classbox] (core_state) at (115mm, -80mm) {\textbf{gameCore\_layout.h}\\\texttt{GameState} (Vòng lặp \& Cờ)};
        \node[classbox] (core_tetro) at (190mm, -80mm) {\textbf{Tetromino Struct}\\\texttt{blocks[4]} \& \texttt{colorID}};
        \node[classbox] (core_point) at (265mm, -80mm) {\textbf{Point Struct}\\\texttt{x, y} (Tọa độ ma trận)};
        \draw[arrow] (core_state) -- (core_tetro);
        \draw[arrow] (core_tetro) -- (core_point);

        \draw[cslate, dashed, line width=1pt] (0, -100mm) -- (295mm, -100mm);

        % ================= LAYER 4: Utilities =================
        \node[layerlabel] at (0, -120mm) {4. Tiện ích Hệ thống};
        \node[classbox, fill=bgbk, draw=cbk] (util_logger) at (152.5mm, -120mm) {\textbf{logger.h}\\\texttt{Logger} (Singleton, \texttt{game.log})};
        \node[classbox, fill=bgbk, draw=cbk] (util_debug) at (227.5mm, -120mm) {\textbf{ctetris\_debug.h}\\(HTTP Macros, WASM/Native)};

        \draw[<->, dashed, line width=1.5pt, draw=cbk] (util_logger) -- (util_debug);
    \end{scope}
  \end{scope}
\end{tikzpicture}
\newpage

\end{document}